% Chapter 3.1: The Multi-Query Bottleneck in Computational Science
% Part III of the book.
% STATUS: rewritten to book standard (Part III draft pass).

\chapter{The Multi-Query Bottleneck in Computational Science}
\label{ch:part3-the-multi-query-bottleneck-in-computationa}
\label{sec:intro_apps}
% -----------------------------------------------------------------------------

Parts~I and~II of this book were written against a physics that stays out of the way. The
domains are the periodic box, the torus, and the unit square; the equations are Darcy flow,
Burgers' equation, and forced two-dimensional Navier--Stokes; the error is a relative $L^{2}$
distance between two fields; and the operator is applied once, to one input, on a grid of a few thousand points. That
austerity was deliberate.
\chapref{ch:part1-operator-approximation-theory-and-universa} and
\chapref{ch:part1-the-fourier-neural-operator-spectral-theor} could state approximation and
convergence guarantees only by holding the geometry simple enough that the solution operator
carries the continuity their hypotheses demand, and
\chapref{ch:part2-canonical-pde-benchmarks-and-reference-wor} could price a training run only by
fixing benchmarks every reader can reproduce\index{multi-query surrogate problem!Part III verification standard|seealso{surrogate verification}}.

Part~III withdraws those concessions one at a time. The physics a laboratory or a clinic actually
pays for arrives with a triangulated surface, a heterogeneous tensor coefficient, a boundary
condition that must hold to tolerance, a horizon of thousands of steps, and --- above all --- a
demand for many solutions of one governing system.
\chapref{ch:part1-introduction-the-paradigm-shift-to-infinit} named that pattern the multi-query
surrogate problem and wrote down its arithmetic; here the same pattern is priced against the
solvers it is meant to displace. Section~\ref{sec:ch31-one-solve} accounts for the cost of a
single classical solve at production resolution and explains why \eqref{eq:ch31-discrete-system}
is a floor rather than an estimate. Section~\ref{sec:ch31-workloads} separates the four workload
families --- posterior sampling, shape optimization, ensemble forecasting, procedural clinical
guidance --- along the axis that matters: not how many queries each issues, but how those queries
relate to one another. Section~\ref{sec:ch31-hypothesis} states the surrogation hypothesis and
repairs its cost accounting; Section~\ref{sec:ch31-fidelity} fixes what counts as evidence for
the five application chapters that follow. No architecture is introduced here and no measurement
is reported for the first time; the chapter is the ledger the rest of the Part is kept
in\index{surrogate verification!standard for Part III}. Notation follows
Table~\ref{tab:notation-master}.

% -----------------------------------------------------------------------------
\section{The Price of One Forward Solve}
\label{sec:ch31-one-solve}
% -----------------------------------------------------------------------------

High-fidelity numerical simulation is the standard instrument of modern engineering, geophysics,
and clinical diagnostics, and its expense is structural rather than accidental. Whether the task
is the compressible Navier--Stokes equations around an aircraft, multi-phase flow in a deep
geological reservoir, or a wavefront crossing a patient-specific ventricle, the classical
workflow discretizes continuous differential operators into large sparse linear or non-linear
algebraic systems,
\begin{equation}
  \mathbf{A}(\mathbf{u}; \mathbf{p})\, \mathbf{u} = \mathbf{f}(\mathbf{p}), \qquad
  \mathbf{u} \in \mathbb{R}^{N}, \quad N \sim 10^{7} \text{--} 10^{9},
  \label{eq:ch31-discrete-system}
\end{equation}
where $\mathbf{p}$ collects the parameters --- boundary data, coefficient fields, geometry ---
and $\mathbf{u}$ stacks the unknown nodal or cell-averaged degrees of freedom
(DOF)\index{degrees of freedom (DOF)!in a production mesh}. Equation
\eqref{eq:ch31-discrete-system} carries three costs, and the exponent on $N$ is only the first.
Because the operator depends on its own solution --- an eddy viscosity $\mu_{t}$ that is a function
of the field being solved, a relative permeability that is a function of the saturation being
tracked --- a converged answer means a nested loop, each outer non-linear step driving an inner
Krylov or multigrid solve of a matrix that has just changed.
Solving such systems requires millions of floating-point operations per time step, and memory
follows $N$ times the number of transported fields, so a mesh at the middle of
\eqref{eq:ch31-discrete-system} is a distributed-memory job, not a workstation one.

Refinement compounds the arithmetic and the memory cost together. In $d$ spatial dimensions a
mesh of spacing $\Delta x$ over a
domain of size $L$ carries $N = (L/\Delta x)^{d}$ unknowns, and an explicit scheme stable at the
Courant--Friedrichs--Lewy (CFL) number $\mathrm{CFL}$ must take
$\Delta t \le \mathrm{CFL}\,\Delta x / |\mathbf{u}|_{\max}$, so that a physical horizon
$T$ fixes the unknown count and the step count,
\begin{equation}
  \label{eq:ch31-refine}
  N(\Delta x) = \Bigl(\frac{L}{\Delta x}\Bigr)^{d}, \qquad
  K(\Delta x) = \frac{T}{\Delta t} \;\ge\; \frac{|\mathbf{u}|_{\max} T}{\mathrm{CFL}\,\Delta x},
\end{equation}
and the arithmetic of one run is their product $N K = \mathcal{O}(N^{1+1/d})$, since $K$
grows like $N^{1/d}$.
Equation \eqref{eq:ch31-refine} is why a refinement study is not a
linear expense: halving
$\Delta x$ in three dimensions multiplies storage by eight and the step count by two. Implicit
treatments buy back the step count and pay in conditioning, since the condition number grows
like $\Delta x^{-2}$ and the iteration count with it unless an optimal preconditioner is carried
along
\cite{leveque2002finite, quarteroni2008numerical}\index{CFL condition!cost of explicit time marching}.

The application literature reports the wall-clock consequences of \eqref{eq:ch31-refine} at the
resolutions industry uses, and those numbers are Part~III's starting ledger. Every one of them
arrives through this Part's working record, and each is audited in the chapter that owns it against
the study named beside it; where that audit finds the study did not print the pairing, the ledger
says so here rather than quietly inheriting a citation. A converged
Reynolds-averaged Navier--Stokes (RANS) volume solve for one road-car geometry, on a body-fitted
mesh with boundary-layer prisms refined to $y^{+} < 1$ and roughly $2 \cdot 10^{7}$ cells, takes
1{,}500--3{,}000 Semi-Implicit Method for Pressure Linked Equations (SIMPLE) or
Pressure-Implicit with Splitting of Operator (PISO) iterations and 15--45 minutes on 64 central
processing unit (CPU) cores. That pairing is the applications literature's: the aerodynamic study
this Part cites, \cite{li2023geometry}, prices its surrogate against a
graphics-processing-unit (GPU) volume solver and reports its own evaluation on a single NVIDIA
V100, as \chapref{ch:part3-external-aerodynamics-and-surface-flows} spells out
\index{Reynolds-averaged Navier--Stokes (RANS)}\index{finite volume method (FVM)!pressure--velocity coupling}.
One global numerical weather prediction (NWP) member integrated with the European Centre for
Medium-Range Weather Forecasts (ECMWF) Integrated Forecasting System (IFS) dynamical core takes
roughly 2 hours on a supercomputer; that elapsed time is inherited, not reported by the
study:
\cite{pathak2022fourcastnet} compares whole ensembles node-to-node and gives no timing for a
single member on one partition\index{numerical weather prediction (NWP)}.
One 30-year injection history in a single realization of the two-phase Darcy system takes about
ten minutes on one dedicated Intel Xeon core in the implicit finite-difference simulator that
generated the labels of \cite{wen2022ufno}, the 10--30 minute range quoted across this Part being
the working envelope carried forward for that class of problem. And a monodomain action-potential
rollout runs to minutes on the secondhand estimate; the workload belongs to
\chapref{ch:part3-cardiac-electrophysiology-and-arrhythmia-d}, whose surrogate is the
physics-informed neural operator (PINO) formulation \cite{li2021physics} applied to the monodomain
equation closed by the Aliev--Panfilov reaction model \cite{aliev1996simple}. No study
this Part cites for that workload prices a data-generating solve against a trained rollout, so
the minutes stand as a carried-forward framing rather than one side of a measured
ratio\index{monodomain equation!cost of a production rollout}.
Where such pairs are on the record, as they are for the reproducible benchmarks,
\chapref{ch:part2-production-deployment-and-high-throughput-} tabulates them --- solver and
surrogate on one configuration --- in Table~\ref{tab:benchmarks_hardware}.

Two cautions belong in the reader's hands before the first application chapter. First, the four
figures above are not comparable: each is a price struck on one machine at one mesh density, and
it is informative only beside the surrogate it was printed against. A
speedup is a property of a configuration, not of an architecture, and Part~III's convention is to
carry baseline and surrogate together or not at
all\index{speedup factor!measured on one configuration|seealso{multi-query surrogate problem}}.
Second, per-solve arithmetic is not the whole of a query: generating a body-fitted mesh is a
per-geometry step with its own failure modes, and neither the ledger nor the cited studies price
it,
so the
unpriced term in \eqref{eq:ch31-discrete-system} is precisely the term that varies when a design
sweep changes shape\index{mesh generation!cost left unreported}. Tomography
inverts the whole picture. Its baseline, filtered backprojection (FBP), is no large sparse solve
but a direct, non-iterative filtering and backprojection, so the computed tomography (CT) chapter
of this Part competes on accuracy under a dose constraint rather than on wall clock: the
difficulty there is ill-posedness, not
arithmetic\index{ill-posed inverse problem}.

% -----------------------------------------------------------------------------
\section{Four Ways to Ask One Equation Many Times}
\label{sec:ch31-workloads}
% -----------------------------------------------------------------------------

Individual forward simulations are routinely executed on high-performance computing (HPC)
clusters, and one simulation justifies no surrogate. What breaks the classical workflow is the
multiplicative structure of real campaigns: a task needing $Q$ solutions of one system costs $Q$
copies of the prices just listed. Table~\ref{tab:ch31-workloads} collects the four families this
Part works through, arranged by the relation among their queries --- chained, neighboring,
perturbed, or independent --- since that relation decides which surrogate machinery earns its
keep.

Bayesian inverse problems are the chained case. Estimating a subsurface permeability field or a
tissue stiffness distribution from scattered measurements replaces the point estimate with a
posterior measure, which under Gaussian measurement error has density
\begin{equation}
  \label{eq:ch31-posterior}
  \frac{\mathrm{d}\mu_{\mathrm{post}}(\theta \mid y)}{\mathrm{d}\mu_{\mathrm{prior}}(\theta)}
  = \frac{1}{Z(\sigma)} \exp\bigl(-\Phi(\theta)\bigr), \qquad
  \Phi(\theta) = \tfrac{1}{2}\,\bigl\| y - \mathcal{G}(\theta) \bigr\|_{\Gamma}^{2}, \qquad
  \Gamma = \sigma^{2} I,
\end{equation}
relative to the prior, where $\mathcal{G}$ is the parameter-to-solution map, $\Phi$ the likelihood
potential, and $\Gamma$ the measurement covariance \cite{stuart2010inverse}. Markov chain Monte
Carlo (MCMC) sampling of \eqref{eq:ch31-posterior} evaluates $\Phi$ once per proposal, so the
forward model runs $\mathcal{O}(10^{5}\text{--}10^{7})$ times --- the count this Part inherits,
stated as an order of magnitude and not a measurement, and in
practice a function of the parameterization and of whether the sampler is
dimension-robust\index{Bayesian inverse problem}\index{Markov chain Monte Carlo (MCMC)!forward evaluations per chain}.
Two features of the chained case matter downstream. Consecutive queries are neighbors in
parameter space, which lets a classical solver warm-start from the previous iterate; a stateless
surrogate discards that advantage. And the deliverable is a measure, not a field: error in $\Phi$
biases the posterior rather than degrading a picture, so accuracy is priced in the likelihood
potential, not in the $L^{2}$ norm.

Shape optimization is the neighboring case, where the query count and the derivative count
multiply. Iterative gradient-based design sweeps explore thousands of geometric perturbations to
minimize drag subject to lift constraints; with $n$ design variables a finite-difference gradient
costs $n+1$ solves per step, which is why production computational fluid dynamics (CFD) adjoints
exist at all\index{computational fluid dynamics (CFD)}. A differentiable surrogate changes that
arithmetic: the adjoint field $\partial \mathcal{L}/\partial v_{\ell+1}$ whose appearance in
\eqref{eq:ch26-mode-gradient} Part~II used to explain a learning-rate schedule is the same object
the design gradient needs, and reverse-mode differentiation through the trained network --- the
autograd engine of \cite{paszke2019pytorch}, whose treatment of complex tensors Part~II also
relies on --- returns $\mathrm{d}\mathcal{G}_\theta/\mathrm{d}\mathbf{p}$ for a small multiple of
one forward pass\index{aerodynamic shape optimization!gradient cost under a surrogate}. The hazard mirrors the benefit:
an optimizer does not evaluate a surrogate where it was trained
but where the surrogate says the objective is smallest, so a systematic bias in the field error
becomes an artificial optimum. Nothing in Part~I bounds that
behavior, because every guarantee there holds on a compact set of inputs
(Assumption~\ref{ass:ch09-compact}) and the point of a design search is to leave the sampled
set\index{surrogate exploitation|seealso{out-of-distribution (OOD) failure}}.

Ensemble forecasting is the perturbed case and the mildest of the four. Characterizing atmospheric chaos
necessitates rolling out 50-member ensembles of global
atmospheric simulations twice daily, each a perturbed initial state of the same system, so the
classical cost is fifty copies of the 2-hour IFS integration above. The queries share one grid and one
model, the best circumstance for the batching argument of
\chapref{ch:part2-production-deployment-and-high-throughput-}: fifty fields ride one batch axis
and the same weights\index{ensemble forecasting}. What is scored, however, is the ensemble's reliability
rather than any single trajectory, so the
fidelity standard moves from field error to skill measures such as the anomaly correlation
coefficient (ACC).

Clinical decision support is the independent-and-novel case and the hardest. Intraoperative guidance
during cardiac ablation or trauma assessment demands sub-second
predictions, which precludes the multi-hour finite element method (FEM) runs that
patient-specific geometry requires; the constraint is a wall clock inside a procedure, not a batch
quota\index{real-time clinical decision support}\index{finite element method (FEM)!patient-scale runs}.
Cardiac electrophysiology approaches that deadline with PINO \cite{li2021physics}, the physics-informed
formulation of the operator problem, which
\chapref{ch:part3-cardiac-electrophysiology-and-arrhythmia-d} trains on the monodomain equation with
Aliev--Panfilov kinetics \cite{aliev1996simple}: the governing residual in the loss is what keeps the
number of expensive labeled solves affordable offline, and it is the trained rollout, not the solve,
that has to fit inside the procedure.
Patient-specific soft-tissue biomechanics is posed against the same deadline with a different kind of map:
there a learned operator runs from an anatomy to the displacement field measured by magnetic
resonance elastography (MRE), which is an anatomy-to-displacement regression rather than a
simulation of a head undergoing impact, and no clinical validation is established for
it\index{soft-tissue biomechanics!elastographic displacement, not impact simulation}.
Hemodynamics sit in the same regime with a different constraint: Velikorodny, Lu, and Dudenkov
\cite{deeponetblood2025} couple a low-fidelity computational-fluid-dynamics surrogate to
high-fidelity in-vitro measurements in a multi-fidelity DeepONet for flow through stenosed
coronary vessels, so that the simulation need not be repeated for each new geometry; the
instantaneous pressure error they report against the in-vitro reference is $1.2 \pm 3.5$
millimeters of mercury (mmHg), and the validation is against phantoms and simulations rather than
prospective patient measurements. Here the scarce
resource is admissible training data and a geometry outside the sampled distribution, not
floating-point operations.

\begin{table}[t]
\centering
\small
\caption{\textbf{The four multi-query workload families of Part~III, organized by the relation
among their queries.} Query counts and per-query costs are those this Part inherited; the
citation names the surrogate study associated with the row, and the row states
explicitly when the study does not itself report the baseline timing beside it. The rows are not
comparable with one another and are not combined anywhere in this book.}
\label{tab:ch31-workloads}
\begin{tabularx}{\textwidth}{@{}l l >{\raggedright\arraybackslash}X
  >{\raggedright\arraybackslash}X@{}}
\toprule
\textbf{Workload} & \textbf{Queries} & \textbf{Relation among queries} & \textbf{Cost per query, as reported} \\
\midrule
Bayesian inversion & $\mathcal{O}(10^{5}\text{--}10^{7})$ & chained; warm-startable & $\approx$10 min per
30-yr realization, one Xeon core \cite{wen2022ufno}; the 10--30 min range is unattributed \\
Shape optimization & thousands & neighboring, driven by gradients & 15--45 min per geometry on 64 CPU cores (unattributed pairing; \cite{li2023geometry} prices a GPU solve) \\
Ensemble NWP & 50 members, twice daily & perturbed initial states, one grid & $\approx$ 2 h per member,
ECMWF IFS (figure with no primary reference; the study times ensembles \cite{pathak2022fourcastnet}) \\
Procedural clinical & one per patient, per step & independent, often novel geometry & minutes per rollout (framing with no primary reference; the associated study reports no solve-versus-rollout pairing for this workload) \cite{li2021physics} \\
\bottomrule
\end{tabularx}
\end{table}

Table~\ref{tab:ch31-workloads} thus says what the query counts conceal: the four families stress
four different properties. Inversion asks whether an error survives
the logarithm in \eqref{eq:ch31-posterior}; optimization asks whether the search stays inside the
training support; ensembles ask whether the spread is honest; clinical use asks what the model may
be wrong about. An architecture shown to serve one family has not been shown to serve the others,
and Part~III reports them separately for that reason.

% -----------------------------------------------------------------------------
\section{The Surrogation Hypothesis and Its Ledger}
\label{sec:ch31-hypothesis}
% -----------------------------------------------------------------------------

The response developed across this Part is the one
\chapref{ch:part1-introduction-the-paradigm-shift-to-infinit} set up in the abstract: instead of
solving \eqref{eq:ch31-discrete-system} once per query, learn the map from data to solution once,
offline, and evaluate the learned map per query. The bookkeeping was written down as
\eqref{eq:ch01-cost-bookkeeping}, where $M$ denoted the query count; from here $Q$ is used,
because $M$ already carries the training-ensemble size of Table~\ref{tab:notation-master} and the
projection-view count of
\chapref{ch:part3-medical-imaging-and-tomographic-reconstruc}. The Part's claim, as this
book received it, follows.

\takeaway{The Operator Surrogation Hypothesis}{Neural operators resolve the multi-query bottleneck by separating computational cost into an expensive offline training phase and a near-instantaneous online evaluation phase. Once trained across a distribution of parameters, the operator predicts continuous solution fields in milliseconds, achieving acceleration factors of $10^{3}\times$ to $10^{5}\times$ over traditional numerical solvers.}

That statement is precise enough to audit, and three items deserve attention. The first is the
cost of the offline phase, which the application literature reports rarely and the headline tables
never. Supervised
training consumes paired samples, and paired samples are produced by running the solver being
replaced:
\begin{equation}
  \label{eq:ch31-ledger}
  c_{\mathrm{train}} \;\approx\; Q_{\mathrm{lab}}\, c_{\mathrm{solve}}(N_{\mathrm{lab}}) \;+\; c_{\mathrm{opt}},
  \qquad
  Q^{\star} \;=\; \frac{c_{\mathrm{train}}}{c_{\mathrm{solve}}(N_{\mathrm{prod}}) - c_{\mathrm{infer}}(N_{\mathrm{prod}})},
\end{equation}
with $Q_{\mathrm{lab}}$ the number of labeled simulations, $c_{\mathrm{opt}}$ the optimizer's own
cost, $N_{\mathrm{lab}}$ and $N_{\mathrm{prod}}$ the labeling and production grid sizes, and
$Q^{\star}$ the break-even query count beyond which the surrogate is cheaper than the solver it
replaces. Every acceleration figure in this Part divides $c_{\mathrm{solve}}$ by
$c_{\mathrm{infer}}$ at one $N_{\mathrm{prod}}$; none divides $C_{\mathrm{sur}}$ by
$C_{\mathrm{sim}}$, and the second quotient is formable only where a study states its label count
and its training time beside its inference time, which this Part does chapter by chapter rather
than here. The form of \eqref{eq:ch31-ledger} also needs one qualification of its own: the
assumption that labels come from the solver being replaced fails exactly once, for the atmospheric
case, where the labels are archived analyses rather than IFS integrations and $c_{\mathrm{train}}$
therefore does not inherit the 2-hour price at all. A headline
speedup therefore licenses nothing on its own: at $Q = 1$ the surrogate is strictly more expensive
than the solver, which is why \eqref{eq:ch01-cost-bookkeeping} was qualified in Part~I with the
instruction to hand a single design check to a classical
solver\index{break-even query count}. The four families of Table~\ref{tab:ch31-workloads} sit on
opposite sides of \eqref{eq:ch31-ledger}. A posterior sweep issues
$\mathcal{O}(10^{5}\text{--}10^{7})$ queries against a labeling set of a few thousand solves, so
its query count exceeds $Q^{\star}$ by two to four orders of magnitude. A fifty-member ensemble
run twice daily accumulates 36{,}500 solves in a year, which clears a labeling set of the same
order and no more. The single reconstruction a clinician orders is one query against a training
cost measured in thousands of solves and never clears $Q^{\star}$ at all, and no accuracy column
changes that arithmetic.

The second item is what makes \eqref{eq:ch31-ledger} tractable, and it is the result of Part~I
that Part~III leans on hardest. If the discretization error of the labels decays with
$N_{\mathrm{lab}}$ at a rate the trained network inherits --- the content of \eqref{eq:ch06-decomp}
and the continuum-limit statement it belongs to --- then $N_{\mathrm{lab}}$ may be chosen far
below $N_{\mathrm{prod}}$, and the prohibitive term $Q_{\mathrm{lab}} c_{\mathrm{solve}}(N_{\mathrm{lab}})$
comes down by the same $N^{1+1/d}$ law that made \eqref{eq:ch31-refine} expensive. Coarse-grid labeling
is not a convenience but the only affordable route to a surrogate for
$N \sim 10^{9}$, and is legitimate precisely to the extent Part~I's discretization invariance holds\index{discretization invariance!as a license for coarse-grid labels}. The companion caveat
is Part~II's \eqref{eq:ch27-budget}: a surrogate's reported error is its error against labels, so what
a campaign may claim is bounded by the accuracy of the grid it was trained on\index{error budget!surrogate against solver labels}.

The third item is the band in the takeaway, whose endpoints are traceable but not all on the same
machine or against the same class of baseline. Roughly $4{,}000\times$ for global weather
\cite{pathak2022fourcastnet} is formed by dividing this Part's own two numbers --- $7{,}200$
seconds of IFS wall-clock against the $1.8$ seconds the working record gives FourCastNet for a
week-long forecast on one NVIDIA A100 --- which \chapref{ch:part3-global-weather-forecasting-and-climate-dyn}
reads as a statement about latency on a device rather than a costing, the study's own comparison
being node-to-node over a whole ensemble. $26{,}000\times$ for car-surface aerodynamics is not a
latency ratio on that device at all, but the cost--accuracy quotient \cite{li2023geometry} forms
against a GPU volume solver, read at a $3\%$ relative error in $C_d$; and $60{,}000\times$ for
30-year $\mathrm{CO}_{2}$ plume migration is an A100 inference ratio against the finite-difference
simulator that generated the labels of \cite{wen2022ufno}, halved to nearer $3\times10^{4}$ in
\chapref{ch:part3-subsurface-multiphase-flow-and-carbon-stor} because saturation and pressure are
returned by two networks. All three land inside $10^{3}\times$ to $10^{5}\times$, which is all the
band claims. One Part~III case falls outside it: tomographic reconstruction. Its baseline is FBP,
the filtering-and-backprojection pipeline that
\chapref{ch:part3-medical-imaging-and-tomographic-reconstruc} derives as its analytic comparator ---
a direct transform returning a slice in milliseconds, so the comparison is not one of arithmetic but
of image quality at lower dose, and a learned reconstruction has no orders of magnitude of runtime
to claim against it. The $\approx 100\times$ printed beside that row in this Part's
ledger is precisely
the kind of quotient \chapref{ch:part3-synthesis-matrix-and-real-world-deployment} now refuses to
carry forward, entering the cell as unestablished, because nothing behind it supplies a pair of
timings and the tomographic chapter of this Part prices its reconstructions on consistency and
dose rather than on wall clock. The band describes solvers whose cost is dominated
by \eqref{eq:ch31-discrete-system}; where the classical method is already a direct transform,
operator learning must justify itself on accuracy, safety, or dose.

% -----------------------------------------------------------------------------
\section{Fidelity Is Scored in the Functional, Not the Field}
\label{sec:ch31-fidelity}
% -----------------------------------------------------------------------------

Every application chapter closes with a fidelity number chosen by what the domain decides rather
than by what the training loss optimized. What a campaign consumes is a functional of the
solution --- an integral, a correlation, a front position, a topological property --- and the
relation between what was minimized and what is reported is
\begin{equation}
  \label{eq:ch31-functional}
  \bigl| \mathcal{J}(\mathcal{G}_{\theta}(a)) - \mathcal{J}(\mathcal{G}(a)) \bigr|
  \;\le\; \bigl\| \mathrm{D}\mathcal{J} \bigr\|_{\mathrm{op}}\,
            \bigl\| \mathcal{G}_{\theta}(a) - \mathcal{G}(a) \bigr\| ,
\end{equation}
where $\mathcal{J}$ is the decision functional and $\|\mathrm{D}\mathcal{J}\|_{\mathrm{op}}$ its
operator norm in the topology of the field error. Equation \eqref{eq:ch31-functional} points the
wrong way to trust: it says a small field error cannot hurt a smooth functional much, while
placing the whole burden on $\|\mathrm{D}\mathcal{J}\|_{\mathrm{op}}$, which is large wherever the
functional reads a boundary, a front, or a peak, and says nothing where the functional is
nondifferentiable exactly where the surrogate errs. Part~II's verification protocol insisted on
the sequence --- relative $L^{2}$ first, then the functional the decision uses
(Section~\ref{sec:ch27-calibration}) --- and the five metrics below realize that second step
domain by domain.

Aerodynamics reports a surface integral. The drag coefficient $C_{d}$, obtained by integrating
pressure and wall shear stress over the body, is the number that sizes an engine, and this Part
reads a surrogate against it at the $3\%$ relative error where \cite{li2023geometry} forms its
cost--accuracy comparison; the tighter agreement of $1.2\%$ quoted in this Part's record is
not a figure that study reports, and \chapref{ch:part3-external-aerodynamics-and-surface-flows}
uses the $3\%$ value. Either way a surrogate may misplace a separation region and still reproduce
the scalar a design office consumes\index{drag coefficient}. Atmospheric science reports a
correlation: the ACC for 500 hPa geopotential height against an analysis, with useful skill taken
as $\mathrm{ACC} > 0.6$ through day 7 in the figures this Part carries forward. The study's own statements are
differently shaped --- ACC above $0.6$ for total column water vapor beyond day eight, and
comparability with the IFS on ACC and root mean squared error only at lead times up to three days
\cite{pathak2022fourcastnet} --- so the horizon is a threshold one agrees to cross, not a property
of the model. ACC measures pattern agreement
relative to a climatology baseline and says nothing about whether a particular storm tracks over
a particular city\index{anomaly correlation coefficient (ACC)}. Carbon storage reports relative field errors,
split by variable --- a mean gas-saturation plume error of $1.6\%$, with a standard deviation near
$1.1$ points, and a mean pressure-buildup error of $0.68\%$ \cite{wen2022ufno}, which are
distributions rather than bounds --- because saturation and pressure are the fields a reservoir engineer reads,
whereas the quantity a licensing body would ask for, where the front meets a fault, is not reported.
Electrophysiology reports the maintenance of reentrant spiral-wave dynamics
and resolution transfer to a finer evaluation mesh rather than a voltage error, because whether an
arrhythmia persists is a topological property, for which
\eqref{eq:ch31-functional} supplies no bound; the surrogate whose behavior is scored that way is the
PINO construction \cite{li2021physics} of \chapref{ch:part3-cardiac-electrophysiology-and-arrhythmia-d},
trained on the monodomain equation with Aliev--Panfilov kinetics \cite{aliev1996simple}, and the
evidence available there is stable reentry on
flat two-dimensional sheets, with wavebreak flagged as collapsing at the extended
horizon\index{reentry!spiral wave}. Tomography is the sharpest instance, because its
three metrics answer three different questions --- peak signal-to-noise ratio (PSNR) in decibels,
the structural similarity index measure (SSIM), and root mean squared error (RMSE) in Hounsfield
units (HU)\index{peak signal-to-noise ratio (PSNR)}\index{structural similarity index measure (SSIM)}
--- and because their meaning is set entirely by the comparator they are read against.
\chapref{ch:part3-medical-imaging-and-tomographic-reconstruc} therefore treats the tomographic
triplet this Part carries as illustrative rather than verified, and no PSNR,
SSIM, or HU figure is reported here as a measurement: an inversion of the raw sparse data with no
regularization imposed is not the comparator a low-dose protocol faces, and against such a baseline
the apparent gain is a statement about the baseline.
RMSE in HU is a field error in the units of the measurement itself, so two reconstructions that
agree under PSNR and SSIM can still differ in
usability --- which is what makes the tomographic chapter of this
Part about ill-posedness rather than about arithmetic.

Three limitations cut across all five and are treated in
\chapref{ch:part3-boundary-condition-enforcement-neural-oper} and
\chapref{ch:part3-synthesis-matrix-and-real-world-deployment} rather than argued here, so this
chapter only registers them. Out of distribution, a surrogate returns a plausible field instead of
a diagnosis of failure, and Part~I's guarantees are silent on inputs outside the compact support
of Assumption~\ref{ass:ch09-compact}; a classical solver at least diverges
visibly\index{out-of-distribution (OOD) failure!a surrogate fails silently}. Conservation is
guaranteed in the average by a least-squares objective and not pointwise, which is the difference
between a good benchmark and a certifiable component\index{conservation laws!not enforced by a least-squares objective}.
And clinical use carries a clearance burden --- prospective, multi-center validation before a
surrogate enters care --- that no accuracy column can discharge\index{regulatory clearance!medical devices and procedural guidance}.

% -----------------------------------------------------------------------------
% Bridge to Chapter 3.2
% -----------------------------------------------------------------------------

The first of Part~III's five domains spends the geometric concession Parts~I and~II reserved. External
aerodynamics is a multi-query workload of the neighboring kind in
Table~\ref{tab:ch31-workloads}: a vehicle surface is an irregular triangulated mesh that changes
with every design perturbation, the governing system is the RANS equations closed by an eddy viscosity that depends on the
solution being solved for, and the query must be answered over a boundary layer resolved to
$y^{+} < 1$ at Reynolds numbers $Re$ of $10^{5}$ and above. Cartesian
spectral methods, which Part~II built on tori, have no native representation for such a surface,
and that is the gap the Geometry-Informed Neural Operator (GINO) construction of
\chapref{ch:part2-graph-neural-operators-gno-and-irregular-g} was designed to
close\index{Geometry-Informed Neural Operator (GINO)!surface-mesh aerodynamics}.
\chapref{ch:part3-external-aerodynamics-and-surface-flows} runs that machinery of point-to-grid
and grid-to-point operators on a family of more than five hundred Ahmed-body variants and reports
surface pressure, wall shear stress, and the drag coefficient built from them --- at a relative
$L^{2}$ surface-pressure error of $8.31\%$ and the $3\%$ drag agreement at which
\cite{li2023geometry} forms its cost--accuracy ratio. The $26{,}000\times$ figure is that ratio,
not a latency quotient on the 64-core machine of this chapter's ledger: \chapref{ch:part3-external-aerodynamics-and-surface-flows}
keeps the $2 \cdot 10^{7}$-cell, 15--45 minute finite volume method (FVM) solve beside the GPU
solve the study actually prices against, which is exactly the distinction this chapter has been at
pains to make --- a speedup is a ratio, not a law. What it adds beyond this ledger is the
accounting \eqref{eq:ch31-functional} demands: which part of a pressure field a drag coefficient
reads, and what a surrogate may be permitted to get wrong.

\vspace{1.5em}
